\section{Experimental Design}
\label{sec:experiments}

What would count as evidence, and what was decided before any test trajectory
ran? This section is design only. The experiments are running, so every cell is
marked \pending{} and every table is a pre-built slot.

\textbf{Splits, and the independent unit.} Seeds come from disjoint reserved
pools of 36 development units, 24 calibration units and 150 test units, six
families by six, four and twenty-five seeds. Each test unit is crossed with the
four information conditions, giving 600 paired replays plus 20 silent-null and 20
false-alert-null controls, so 640 held-out rollouts per sparse arm, with every
threshold and mapping frozen on development and calibration data. \textbf{The independent unit is the seed, not the rollout.} The four
information cells of one seed are paired replays of one exogenous trajectory, so
treating them as four samples would inflate the effective sample size fourfold
and narrow every interval by roughly half, invisibly. The evaluation engine is
therefore registered to refuse any confirmatory statistic computed from a frame
with more than one row per pair of independent unit and arm, and to raise rather
than warn. Family tail estimates stay exploratory even at twenty-five seeds, and
every language-level claim bootstraps jointly over seed and template identity.

\textbf{Endpoints.} The primary reward endpoint is cumulative undiscounted profit
per episode and the primary safety endpoint is total lost-sales units per episode,
both macro-averaged with equal weight across the six family strata and one null
stratum and across the four conditions within a family. Fixing endpoint and
weighting in advance prevents a later choice among raw profit, a ratio and a
favorable mixture.

\textbf{The three contrasts.} Three contrasts and two endpoints give six tests,
Holm-corrected as one family, and Table~\ref{tab:ladder} holds their slots. C1 is arm 10 against arm 8: does verification earn
its place, or is immediate activation enough? C2 is arm 10 against arm 6: does
the typed hypothesis beat ephemeral numerical parameters at matched calls? C3 is
arm 10 against arm 11: does it beat the strongest published method? C1 and C2 are
mechanism contrasts and hold the trigger time, the information, the controller,
the model, the decoding configuration and the realized call budget fixed, so the
treatment is a single factor. C3 is not, and we do not present it as one, because
arm 11 differs in timing, content and trust at once. Arm 2 is the registered
secondary comparator for whether language is needed at all. Primary inference is
two-sided paired randomization on per-seed differences after condition cells are
aggregated within their seed, with seed-level cluster-bootstrap intervals and
Wilcoxon tests as a secondary check, and every analysis not listed here is
exploratory.

\textbf{Three further reporting rules.} False activation and wrong-family
activation occupy separate columns and are never merged, because the first is
covered by Theorem~\ref{thm:activation} and the second has no theorem-backed
bound. Forty main-table null rollouts cannot estimate a five percent error rate,
so a separate bank of 1{,}300 episodes runs the complete pipeline: 1{,}000
well-specified registered-model nulls across five strata, and 300 that violate
those assumptions, reported separately and \emph{never} pooled with them. And
because an earlier call opportunity must not be mistaken for evidence that a
message meant something, every semantic-content comparison fixes the alert
timestamp, the trigger trace hash and the call opportunity across five variants,
while the early-warning comparison is a \emph{different estimand}. Efficiency
reporting covers calls, repairs, rejections, tokens, latency quantiles, dated
dollar cost and actions per accepted call, with profit-against-cost frontiers on
\emph{realized} budgets. One always-online sweep is
$720 \times 50 + 600 \times 47 = 64{,}200$ calls, roughly \$108 to \$217 at the
dated primary price.
\pending{activation ablation, null-audit calibration and fragility, content
against timing, efficiency and frontiers, operational and hypothesis metrics,
traced cases}
